\documentclass[11pt]{article}
\usepackage[margin=1in]{geometry}
\usepackage{amsmath,amssymb,amsthm}
\usepackage{graphicx}
\usepackage[hidelinks]{hyperref}
\setlength{\emergencystretch}{2em}

\newtheorem{theorem}{Theorem}
\newtheorem{lemma}{Lemma}
\newtheorem{remark}{Remark}

\newcommand{\R}{\mathbb R}
\newcommand{\N}{\mathbb N}
\newcommand{\Hilb}{\mathcal H}
\newcommand{\dist}{\operatorname{dist}}

\title{A Logarithmic Hilbert Compression Gap Dichotomy}
\author{Run \texttt{fa\_banach\_001}, lane 19}
\date{June 19, 2026}

\begin{document}
\maketitle

\section*{Source Problem}

The source paper is Goulnara Arzhantseva, Cornelia Drutu, and Mark Sapir,
\emph{Compression functions of uniform embeddings of groups into Hilbert and
Banach spaces}, arXiv:math/0612378.  Its Question 5.22 asks:

\begin{center}
\includegraphics[width=0.95\textwidth]{figures/open_problem_crop.png}
\end{center}

\noindent
In the terminology of Definition 3.3 of the source paper, a pair \((f,g)\)
is a Hilbert compression gap for a metric space \(X\) if there is an
\(f\)-embedding of \(X\) into Hilbert space and every Hilbert lower compression
function \(\rho\) satisfies \(\rho\ll g\).  The relation \(a\ll b\) is the
usual large-scale comparison: \(a(t)\le C b(Ct)+C\).

\section*{Answer}

The literal formulation has a small but important convention issue: a group
which does not uniformly embed into Hilbert space cannot have any proper
Hilbert compression gap under the strict reading of Definition 3.3.  The
source paper itself uses logarithmically degraded lower functions in Theorem
5.5, so the result below is stated in that same large-scale compression-gap
convention and flags the stricter proper-lower-function reading explicitly.

\begin{theorem}[logarithmic annular compression envelope]
Let \(G\) be an infinite finitely generated group with a word metric.  If
\(G\) admits a uniform, equivalently coarse, embedding into Hilbert space, then
there is a proper nondecreasing function \(f:\R_+\to\R_+\) such that \(G\) has,
in the source paper's large-scale sense, a Hilbert space compression gap of
the form
\[
        \left(\frac{f}{\log(x+2)}, f\right)
\]
The same construction gives a strict Definition 3.3 compression gap whenever
this lower function \(f/\log(x+2)\) is proper.  Conversely, if a group has a
strict proper Hilbert compression gap, then it is Hilbert coarsely embeddable.
\end{theorem}

\section*{Proof Intuition}

For each dyadic annulus \(2^k\le d(x,y)<2^{k+1}\), measure the largest
separation that a single \(1\)-Lipschitz Hilbert map can force on every pair
in that annulus.  Call this number \(H_k\).  If the group coarsely embeds,
then these annular optima go to infinity.

The upper compression function should not be \(H_k\) itself, since annular
optima may oscillate.  The tail infimum
\[
        F_k=\inf_{j\ge k} H_j
\]
is monotone and still tends to infinity.  It is a universal upper bound, up to
large-scale comparison, because any lower control at a given scale must also
be valid at all larger dyadic scales.

For the lower bound, choose almost optimal maps on all annuli and put them in
a Hilbert direct sum with coefficient \(1/k\) on the \(k\)-th annulus.  The
coefficients are square summable, so the sum is Lipschitz, while on the
\(k\)-th annulus the \(k\)-th coordinate alone gives the lower bound
\[
        \frac{H_k}{k}\ge \frac{F_k}{k}\asymp \frac{f}{\log x}.
\]

\section*{Proof}

We use the source paper's convention that a \(\rho\)-embedding into Hilbert
space is a Lipschitz map \(\phi:G\to\Hilb\) such that
\[
        \|\phi(x)-\phi(y)\|\ge \rho(d(x,y))
\]
for all \(x,y\in G\), where \(\rho\) is proper and nondecreasing.

Assume that \(G\) uniformly embeds into Hilbert space.  Fix a word
metric \(d\).  For each integer \(k\ge1\), define
\[
 A_k=\{(x,y)\in G^2:2^k\le d(x,y)<2^{k+1}\}.
\]
Let \(H_k\) be the supremum, over all \(1\)-Lipschitz maps
\(\phi:G\to\Hilb\), of
\[
        \inf_{(x,y)\in A_k}\|\phi(x)-\phi(y)\|.
\]
The numbers \(H_k\) are finite, since every \(1\)-Lipschitz map separates
pairs in \(A_k\) by at most \(2^{k+1}\).

Because \(G\) uniformly embeds into Hilbert space, after rescaling there is a
\(1\)-Lipschitz map \(\psi:G\to\Hilb\) and a proper nondecreasing function
\(\eta\) such that
\[
        \|\psi(x)-\psi(y)\|\ge \eta(d(x,y)).
\]
Thus \(H_k\ge \eta(2^k)\), and \(H_k\to\infty\).

Set
\[
        F_k=\inf_{j\ge k} H_j .
\]
Then \(F_k\) is nondecreasing and \(F_k\to\infty\).  Define a proper
nondecreasing function \(f:\R_+\to\R_+\) by setting, for large \(t\),
\[
        f(t)=F_k \qquad\text{when } 2^k\le t<2^{k+1},
\]
and extending arbitrarily on bounded intervals.

\subsection*{Universal upper bound}

Let \(\rho\) be the lower control of any Hilbert \(\rho\)-embedding
\(\phi:G\to\Hilb\).  Suppose \(\phi\) is \(L\)-Lipschitz.  Then \(L^{-1}\phi\)
is \(1\)-Lipschitz, so for every \(j\),
\[
        L^{-1}\rho(2^j)\le H_j.
\]
If \(2^k\le t<2^{k+1}\), then for all \(j\ge k+1\),
\[
        \rho(t)\le \rho(2^j)\le L H_j.
\]
Taking the infimum over \(j\ge k+1\) gives
\[
        \rho(t)\le L F_{k+1}\le L f(2t).
\]
Therefore \(\rho\ll f\).  Since \(\rho\) was arbitrary, \(f\) is an upper
bound for a Hilbert compression gap of \(G\).

\subsection*{Lower bound with one logarithmic loss}

For each \(k\), choose a \(1\)-Lipschitz map
\(\phi_k:G\to\Hilb_k\) such that
\[
        \inf_{(x,y)\in A_k}\|\phi_k(x)-\phi_k(y)\|\ge \frac12 H_k.
\]
Fix \(e\in G\), replace \(\phi_k\) by \(\phi_k-\phi_k(e)\), and define
\[
        \Phi(x)=\bigoplus_{k=1}^\infty \frac{1}{k}\phi_k(x).
\]
For each \(x\), this vector is well-defined because
\(\|\phi_k(x)\|\le d(x,e)\) and \(\sum k^{-2}<\infty\).  Moreover,
\[
        \|\Phi(x)-\Phi(y)\|^2
        \le d(x,y)^2 \sum_{k=1}^\infty \frac{1}{k^2},
\]
so \(\Phi\) is Lipschitz.

If \(d(x,y)\in[2^k,2^{k+1})\), then the \(k\)-th coordinate gives
\[
        \|\Phi(x)-\Phi(y)\|
        \ge \frac{1}{k}\|\phi_k(x)-\phi_k(y)\|
        \ge \frac{H_k}{2k}
        \ge \frac{F_k}{2k}.
\]
Since \(k\asymp \log(d(x,y)+2)\) and \(f(d(x,y))=F_k\) on this annulus,
\(\Phi\), after a harmless rescaling, is an
\[
        \frac{f(x)}{\log(x+2)}
\]
embedding.  Hence \((f/\log(x+2),f)\) is a Hilbert compression gap.

The functions \(\log x\) and \(\log(x+2)\) are equivalent for the source
paper's large-scale comparison, so this is the requested form.

Finally, if one insists that the lower function in a compression gap be
proper, then any group with such a gap is Hilbert coarsely embeddable by
definition.  Thus non-Hilbert-embeddable finitely generated groups, whose
existence is recalled in the source paper's introduction via Gromov's
construction, are counterexamples to the literal all-groups reading under that
strict convention.
\qed

\section*{Verification and Novelty Search}

The proof is purely metric and uses only Hilbert direct sums and dyadic
annuli.  The bounded novelty check used the following searches:
\[
\begin{gathered}
\texttt{Hilbert space compression gap f/log},\\
\texttt{compression gap coarsely embeds Hilbert log},\\
\texttt{optimal compression function coarse embedding Hilbert log},\\
\texttt{arXiv:math/0612378 compression gap}.
\end{gathered}
\]
Semantic Scholar metadata for arXiv:math/0612378 showed 35 citations; titles
and abstracts containing compression, gap, amenability, wreath products, and
Hilbert compression were inspected, with no direct later answer to Question
5.22 found.  The result should nevertheless be reviewed for possible folklore
status, since the dyadic direct-sum argument is natural.

\section*{References}

\begin{thebibliography}{9}

\bibitem{ADS}
Goulnara Arzhantseva, Cornelia Drutu, and Mark Sapir,
\emph{Compression functions of uniform embeddings of groups into Hilbert and
Banach spaces}, arXiv:math/0612378.

\bibitem{Gromov}
Mikhail Gromov,
\emph{Random walk in random groups}, Geom. Funct. Anal. 13 (2003),
73--146.

\end{thebibliography}

\end{document}
